\section{Minimal singularities on blowups}\label{sec:blowups}

On a blowup, the background itself can force an exceptional divisor into every positive current. The coefficient of this divisor is then part of the minimal ordinary type, while the rest of the potential descends to the target of the blowup, so the equation reduces to one on the target. The density changes by the Jacobian of the blowup, which determines the loss of integrability.

In this section the roles of the letters differ from those in \cref{sec:prescribed}: $X$ is a blowup and $Y$ its target. Let $(Y,h)$ be a connected compact Hermitian $n$-fold, $n\geq2$, let $Z\subset Y$ be a connected smooth center of codimension $r\geq2$, and write $\pi\colon X=\operatorname{Bl}_ZY\to Y$ for the blowup, with exceptional divisor $E$. Let $s_E$ be the canonical section of $\mathcal O(E)$, choose a smooth Hermitian metric on $\mathcal O(E)$ with curvature form $\Theta_E$, so that $\ddc\log|s_E|^2=[E]-\Theta_E$, put $L=\log|s_E|^2$, and fix $c>0$. The smooth background
\begin{equation}\label{eq:blowup-background}
 \theta=\pi^*h+c\Theta_E
\end{equation}
can be indefinite, and it is nonclosed when $h$ is nonclosed. Positivity is imposed on $\theta+\ddc u$, not on $\theta$.

\begin{proposition}\label{prop:blowup-representation}\index{Minimal singularity!exceptional poles}
Every $u\in\PSH(X,\theta)$ has, up to an additive constant, a unique $h$-plurisubharmonic target potential $q$ such that
\begin{equation}\label{eq:blowup-representation}
 u=cL+q\circ\pi+C,
 \qquad
 \theta+\ddc u=\pi^*(h+\ddc q)+c[E].
\end{equation}
The potentials with minimal ordinary singularities are exactly those of the form $cL+O(1)$, and $u$ is one of them if and only if $q$ is bounded. There is a strict $\theta$-plurisubharmonic barrier with logarithmic coefficient $c+\delta$ for some $\delta>0$.
\end{proposition}

\begin{proof}
Put $T=\theta+\ddc u$ and $q=\pi_*(u-cL)$ as a scalar distribution. The distribution $q$ is locally integrable: off $Z$ it is the function $(u-cL)\circ\pi^{-1}$, and a smooth target volume pulls back with bounded Jacobian, so its absolute integral is controlled by that of $u-cL$ on $X$. Since $\pi_*[E]=0$,
\[
 S=\pi_*T=h+\ddc q\geq0.
\]
A smooth local strictly plurisubharmonic shift gives a quasi-plurisubharmonic representative of $q$. Its pullback $q\circ\pi$ is not identically $-\infty$ and satisfies $\pi^*(h+\ddc q)\geq0$, by \cref{lem:smooth-disk-criterion}: since $\pi$ is dominant, the composition of a local plurisubharmonic function with $\pi$ is again plurisubharmonic and not identically $-\infty$.

Off $E$, $T=\pi^*S$. Hence the difference $R=T-\pi^*S$ is an order-zero current supported on $E$, and it is closed: $\mathrm{d}T$ and $\mathrm{d}(\pi^*S)$ both equal $\pi^*\mathrm{d} h$. In a chart where $E=\{z=0\}$, order zero gives $zR=\bar zR=0$. Closedness then gives $\mathrm{d}z\wedge R=\mathrm{d}\bar z\wedge R=0$. Thus only its normal $(1,1)$ coefficient remains, a measure on $E$; the tangential equations in $\mathrm{d}R=0$ make that measure a constant multiple of integration on $E$. Connectedness of $E$ identifies one constant $a$ globally, so $R=a[E]$. Since $R=c\Theta_E+\ddc(u-q\circ\pi)$, its cohomology class is $c[E]$. On a line in an exceptional fiber, $\mathcal O(E)$ has degree $-1$, so $[E]$ is nonzero and $a=c$. The difference between the two scalar potentials in \eqref{eq:blowup-representation} has $\ddc$ zero. Local elliptic regularity makes it pluriharmonic; compactness and connectedness make it constant.

Since $h>0$, $q=0$ is admissible, and every $h$-plurisubharmonic $q$ is bounded above on compact $Y$. Hence $cL$ is minimal and two-sided bounded comparison with it is exactly boundedness of $q$. For the barrier choose another divisor metric, with logarithmic norm $L'=L+\sigma$ and curvature $\Theta'_E=\Theta_E-\ddc\sigma$ for a smooth real function $\sigma$, so that $-\Theta'_E$ is positive on the vertical tangent bundle of $E\to Z$, using the tautological $\mathcal O(-1)$. In adapted coordinates
\begin{equation}\label{eq:blowup-chart}
 \pi(y,z,w)=(y,z,zw),
\end{equation}
$\pi^*h$ controls the nonvertical directions and $|z|^2|\mathrm{d} w|^2$. For small $\delta>0$, the vertical positivity of $-\delta\Theta'_E$ dominates its mixed errors by completing the square; away from $E$, compactness allows the same $\delta$. Thus $\pi^*h-\delta\Theta'_E>0$. For the fixed background $\theta$ in \eqref{eq:blowup-background}, take
\[
 \rho=(c+\delta)L'-c\sigma=(c+\delta)L+\delta\sigma.
\]
Then $\theta+\ddc\rho=(c+\delta)[E]+\pi^*h-\delta\Theta'_E$ dominates a positive Hermitian form. The smooth correction $\delta\sigma$ preserves the stated logarithmic coefficient.
\end{proof}

\Needspace{18\baselineskip}
\begin{theorem}\label{thm:minimal-blowup-data}\index{Blowup!density integrability}
Fix a smooth positive volume form $\mu$ on $X$ and a nonzero nonnegative $f\in L^p(X,\mu)$, $p>1$. There exist $\kappa>0$ and a potential $u$ with minimal ordinary singularities such that
\begin{equation}\label{eq:minimal-blowup-equation}
 \mathrm{MA}^{\mathrm{ext}}_\theta(u)=\kappa f\mu,
 \qquad u=cL+q\circ\pi+C,
\end{equation}
where $q$ is H\"older on $Y$ and $u$ is normalized by $\sup_Xu=0$. The divisor coefficient is exactly $c$. In \eqref{eq:blowup-chart}, $u-cL=q(y,0)+C+O(|z|^\alpha)$ on smaller charts for some $\alpha>0$, so its limit is constant along each exceptional fiber.

If $\pi_*(f\mu)=Fh^n$, then $F\in L^s(Y,h^n)$ for every
\begin{equation}\label{eq:pushforward-exponent}
 1<s<\frac{rp}{(r-1)p+1}.
\end{equation}
For a point center and data bounded above and below by positive constants near $E$, the density $F$ is $L^s$-integrable near $Z$ exactly for $s<n/(n-1)$, and for finite $p$ no strictly larger universal exponent than $np/((n-1)p+1)$ can be guaranteed. If $h$ is K\"ahler, the class $[\theta]$ is big and non-nef, has volume $\int_Yh^n$, and $\kappa=\int_Yh^n/\int_Xf\mu$; the normalized solution is unique among all normalized solutions of \eqref{eq:minimal-blowup-equation} with minimal ordinary singularities. Under $\ddc(h^j)=0$, $1\leq j\leq n-1$, the scalar of the constructed solution has the same integral formula.
\end{theorem}

\begin{proof}
\emph{Integrability and existence.} In \eqref{eq:blowup-chart} the complex Jacobian is $z^{r-1}$. Put $J=\pi^*h^n/\mu$, so $J\asymp|z|^{2(r-1)}$ on nested charts and $F\circ\pi=f/J$ off $E$. For $1<s<p$, change of variables and H\"older's inequality give
\begin{align}\label{eq:pushforward-holder}
 \int_YF^s h^n
 &=\int_X f^sJ^{1-s}\mu\\
 &\leq\|f\|_{L^p(\mu)}^s
 \left(\int_XJ^{-(s-1)p/(p-s)}\mu\right)^{(p-s)/p}.
\end{align}
The normal complex integral is finite precisely if $(r-1)(s-1)p/(p-s)<1$, which is \eqref{eq:pushforward-exponent}. \Cref{thm:kn-continuous,thm:hermitian-holder-elementary} on $(Y,h)$ supply $q$ and $\kappa$ with $(h+\ddc q)^n=\kappa Fh^n$. Use \eqref{eq:blowup-representation} off $E$ and \cref{thm:arbitrary-extended-ma} across it. H\"older continuity gives $q(\pi(y,z,w))=q(y,0)+O(|z|^\alpha)$ uniformly on smaller charts.

\emph{Sharpness of the exponent.} For $f$ bounded above and below by positive constants near the exceptional divisor of a point center, $F$ is comparable to $\operatorname{dist}_h(\cdot,Z)^{-2(n-1)}$ near $Z$, giving the strict threshold. More generally, $f\asymp|s_E|^{-2a}$ with $0\leq a<1/p$ belongs to $L^p$ upstairs, while its pushforward has threshold $s<n/(n-1+a)$. Letting $a\uparrow1/p$ proves the stated limitation; it makes no finite-$p$ endpoint assertion.

\emph{The K\"ahler case.} When $h$ is K\"ahler, the strict barrier current $(c+\delta)[E]+\pi^*h-\delta\Theta'_E$ constructed in \cref{prop:blowup-representation} makes $[\theta]$ big. Its degree on an exceptional-fiber line is $-c$, so it is not nef. The minimal current $c[E]+\pi^*h$ has non-pluripolar top mass $\int_Yh^n$, and the minimal mass theorem identifies this with $\operatorname{vol}([\theta])$ as recalled in \cref{sec:history-closed-loci}. Full-mass uniqueness and mass integration fix the scalar and normalized current as recalled in \cref{sec:history-regularity}. Another normalized solution with minimal ordinary singularities, possibly discontinuous, still has bounded target potential by \cref{prop:blowup-representation}. It is bounded on every relatively compact subset of $X\setminus E$, so its deep truncations coincide with it there, and locality in \cref{thm:arbitrary-extended-ma} gives the bounded target equation off $Z$. Both target measures give $Z$ zero mass, so closed uniqueness applies globally. If only the power conditions hold, \cref{thm:mass-invariance-characterization} fixes the scalar of the constructed solution. Quantitative H\"older bounds depend on $\kappa F$, not on an unnormalized $F$ alone.
\end{proof}

For the prescribed pullback measure $\pi^*(h^n)$, choose $q=0$ and $\kappa=1$. The current is $c[E]+\pi^*h$. Its smooth part $\pi^*h$ is degenerate on $E$ exactly in the directions tangent to the exceptional fibers, and its density relative to a smooth volume form is $|z|^{2(r-1)}$ times a smooth positive factor. On $X\setminus E$, where the solution $u=cL$ is smooth and strictly $\theta$-plurisubharmonic, the metric $\theta+\ddc u=\pi^*h$ is isometric to $(Y\setminus Z,h)$; its metric completion is $Y$, so each exceptional fiber collapses to a point of $Z$. The distance to $E$ is comparable to $|z|$ on bounded adapted charts. Neither the density order nor this distance exponent depends on $c$.

\begin{example}\label{ex:surface-blowup}
If $n=2$ and $Z$ is a point on a compact K\"ahler surface, $[\theta]=\pi^*[h]+c[E]$ has Zariski positive part $\pi^*[h]$ and negative part $cE$: $E^2=-1$ and $\pi^*[h]\cdot E=0$. Its volume is $\int_Yh^2$, although it is non-nef. The solution with $\kappa=1$ for the measure $\pi^*(h^2)$ has logarithmic coefficient exactly $c$, its density vanishes to order $|z|^2$ along $E$, and the completion of its metric on $X\setminus E$ collapses $E$ to the point $Z$. The stronger barrier from \cref{prop:blowup-representation} has coefficient $c+\delta$.
\end{example}

\Needspace{30\baselineskip}
\begin{theorem}\label{thm:iterated-blowups}\index{Blowup!iterated equations}
Let $\pi\colon Y_m\to Y_0$ be a finite nonempty sequence of blowups $\pi_j\colon Y_j\to Y_{j-1}$ of connected smooth centers of codimensions $r_j\geq2$ over a compact Hermitian $(Y_0,h_0)$. Let $E_j$ be the exceptional divisor of $\pi_j$, let $\pi_{>j}\colon Y_m\to Y_j$ denote the composition of the subsequent blowups, choose $a_j>0$, and set
\[
 \Delta=\sum_j a_j\pi_{>j}^*E_j,
 \qquad \ddc L_\Delta=[\Delta]-\Theta_\Delta,
 \qquad \theta=\pi^*h_0+\Theta_\Delta,
\]
where $L_\Delta=\sum_ja_j\log|s_j|^2$, with $s_j$ the canonical section of $\pi_{>j}^*\mathcal O(E_j)$ for a smooth Hermitian metric, and $\Theta_\Delta$ is the corresponding combination of curvature forms. Every $\theta$-positive current is $\pi^*(h_0+\ddc q)+[\Delta]$ and every potential is $L_\Delta+q\circ\pi+C$, with $q\in\PSH(Y_0,h_0)$. Minimal ordinary type is exactly $L_\Delta+O(1)$, equivalent to bounded $q$. For potentials of this type, every prime exceptional coefficient is the one prescribed by the effective total transform $\Delta$. There is a strict barrier with divisor $\Delta+B$ for an effective exceptional $B$ positive on every exceptional component. For nonzero nonnegative upstairs $L^p$ data, $p>1$, a minimal solution has H\"older $q$, with pushforward density in every $L^s$ satisfying
\begin{equation}\label{eq:iterated-exponent}
 1<s<\frac{Rp}{(R-1)p+1},\qquad R=\prod_jr_j.
\end{equation}
For K\"ahler $h_0$ the mass, scalar, and closed full-mass uniqueness are those of the target; the power conditions on $h_0$ fix the scalar for the constructed nonclosed solution. For the pullback measure $\pi^*(h_0^n)$, the metric $\pi^*h_0$ on the complement of the exceptional set has metric completion $Y_0$, each fiber collapsing to a point, and the zero divisor of the density is that of the squared modulus of the holomorphic Jacobian, independently of the $a_j$.
\end{theorem}

\begin{proof}
Apply \cref{prop:blowup-representation} successively, allowing an intermediate smooth background to be indefinite. The difference between each current and the pullback of its pushforward is closed of order zero, and exceptional-fiber degree fixes its coefficient. Pullback of earlier divisor currents produces the total transforms, proving the representation without an SNC requirement. The initial $q=0$ is admissible; upper bounds on compact targets prove minimality and the bounded-$q$ criterion. At SNC charts this is the usual sum of logarithms.

At each stage choose an auxiliary divisor metric as in \cref{prop:blowup-representation}, with curvature $\Theta'_{E_j}$, so that $h_j=\pi_j^*h_{j-1}-\delta_j\Theta'_{E_j}>0$ for a sufficiently small $\delta_j>0$. Expanding yields $h_m=\pi^*h_0-\Theta_B$ for $B=\sum_j\delta_j\pi_{>j}^*E_j$, effective with a positive coefficient on every exceptional component. If $\ddc L_B=[B]-\Theta_B$ uses these auxiliary metrics, then the fixed background satisfies $\theta+\ddc(L_\Delta+L_B)=h_m+[\Delta+B]$, which is the strict barrier current. The Jacobian calculation and H\"older estimate apply stagewise. If $s=rp/((r-1)p+1)$, then $1-1/s=(1-1/p)/r$; iteration over the stages gives \eqref{eq:iterated-exponent}, with strict inequality. The target H\"older equation, mass and uniqueness are obtained as in \cref{thm:minimal-blowup-data}. For the pullback measure take $q=0$ and use the target path comparison away from the exceptional set. The Jacobian fixes the density, while $[\Delta]$ fixes the divisor current; these are independent pieces of the equation.
\end{proof}
